% Nouveau EGA V (cote 25): archive cote 25, site page 80 (PDF page 81).
% typescript. Draft compared: none.
% AI-transcribed by Claude Opus 5.5 (claude-opus-5-5, adaptive thinking, effort high), Message Batches API, batch ega25b (msgbatch_01A7n5YuXubUPNK1zsepi8RG), 2026-11-13; not yet checked by eye.
\scanpage{25/80}
\mspage{\hand{75}}

\margin{|}m, en notant qu'à changement de corps de base près, $\underline{X}^{(m)}_{\hand{\unsure{\eta|\xi}}}$ s'obtient en
prenant une L' de codimension m-1 passant par x, une H de codi\xout{me}nsion
1 passant par x, L' et H étant génériques indépendants pour ces proprié\xout{x}
tés (i.e. en termes orthodoxes, on se place au point générique du \xout{xx}
schéma des couples (L',H) \xout{xxxxxx} ) et en prenant la section linéaire
de X par $L'\cap H$ : on peut commencer par prendre \xout{\unsure{l}}a section par H, qui
est lisse par ce qui précède, puis par L', qui est lisse par l'hypothèse
de récurrence. Ce genre de raisonnement, déjà utilisé\xout{e} pour généraliser
2.6. p.ex. aux sections linéaires de codimension m queconque, mérite
d'être explicité une bonne fois en général, pour pouvoir y référer sans
entrer \xout{dans} chaque fois dans les détails un peu lourds d'une présentati-%
on en forme.

Reste à prouver le corollaire annoncé de 8.8. dans le cas m=1 .
Lorsque \xout{dim} codim T $\geqslant$ 2, comme d'autre part l'hyperplan \inserted{Q} de $\overset{\hand{\vee}}{P}$ des
$H_{\hand{\xi}}$ tels que $x\in H_{\hand{\xi}}$ est de codimension 1, son point générique \inserted{$\eta$} ne peut
être élement de T, et on gagne. Dans le cas codim T = 1, comme T est
irréductible, on ne peut avoir $\hand{\eta}\in T$ que si Q=T ,\xout{xxxxxxxxxxxxx}
i.e. en termes géométriques \xout{xxxx} \inserted[typed]{(supposant k alg clos, ce qui est loisible)} pour tout $z\in X(\unsure{k})$\xout{xxx}l'espace tangent
de X en z (ou plutôt son image par $f'_{z}$ ) passe\xout{x} par x. Montrons que
ceci ne peut se présenter lorsque $\underline{X}^{\mathrm{sing}}\hand{$\longrightarrow$}T$ est gén étale, i.e.
lorsqu'on est sous les conditions de 8.8., sauf si f(X) = \hand{\{}x\hand{\}} donc X
de dimension 0. En effet, 8.7. c) (qui exprime essentiellement la
symétrie \xout{xxxxx} dans la relation entre X et sa \hand{``}duale\hand{''} \inserted{T}\hand{)} implique alors que
pour presque tout point $z\in X(k)$, f(z) est orthogonal à l'espace tangent
à T en un certain point,\xout{xx} donc (puisque T=Q) \inserted{orthogonal} à Q, d'où f(z)=x, d'où
f(X) =\hand{\{}x\hand{\}} . Cela achève la démonstration de nos remords, donc de 13.2.

\margin{80}
